\section{Limites em espaços métricos}

Como na seção anterior, ao longo dessa seção, a menos que seja explicitamente dito o contrário, $(M, d)$ e $(N, d')$ são espaços métricos arbitrários.

Na seção anterior, vimos que a continuidade de uma função é trivialmente verdadeira nos pontos isolados de seu domínio.

Nos pontos não isolados de seu domínio, ou seja, nos pontos de seu domínio que são pontos de acumulação, a continuidade de uma função pode ser ou não ser verdadeira, dependendo da função.

Uma ferramenta muito útil para testar a continuidade de uma função nesses pontos é a noção de limite.
Tal noção também dá informações sobre o comportamento de uma função nas imediações de pontos de acumulação de seu domínio que estejam em um espaço ambiente, mas que não pertençam ao domínio.

Intuitivamente, $L$ é limite de uma função $f$ no ponto $p$ se, para qualquer margem de erro em torno de $L$, existe uma margem de segurança em torno de $p$ tal que, para todo ponto do domínio de $f$ dentro dessa margem de segurança e distinto de $p$, a imagem desse ponto por $f$ está dentro da margem de erro em torno de $L$.
O ponto $p$ é excluído dessa análise pois a noção de limite pergunta o que ocorre perto do ponto $p$, mas não exatamente no ponto $p$.

\begin{definition}\index{limite}\index{limite!definição}
    Sejam
    \begin{itemize}
        \item $M, N$ espaços métricos,
        \item $A\subseteq M$,
        \item $f: A \to N$,
        \item $p \in M$ um ponto de acumulação de $A$,
        \item $L \in N$.
    \end{itemize}
    Dizemos que $L$ é limite de $f$ em $p$ se, para todo $\epsilon > 0$, existe $\delta > 0$ tal que, para todo $x \in A\setminus\{p\}$,
    \[
        \text{se } d(x, p) < \delta, \text{ então } d'(f(x), L) < \epsilon.
    \]
\end{definition}


Conforme veremos mais adiante, nem sempre uma função possui limites em pontos dados.
Porém, limites, quando existem, são únicos.


A Figura~\ref{fig:limiteUnico} resume a escolha de $\delta$ na prova da Proposição~\ref{prop:limiteUnico}.

\begin{proposition}\index{limite!unicidade}\label{prop:limiteUnico}
    Sejam
    \begin{itemize}
        \item $M, N$ espaços métricos,
        \item $A\subseteq M$,
        \item $f: A \to N$,
        \item $p \in M$ um ponto de acumulação de $A$,
        \item $L_1, L_2 \in N$.
    \end{itemize}

    Se $L_1$ e $L_2$ são limites de $f$ em $p$, então $L_1 = L_2$.
\end{proposition}

\begin{proof}
    Suponha por absurdo que $L_1\neq L_2$. Então $d(L_1, L_2) > 0$.

    Seja $R=\frac{d(L_1, L_2)}{2} > 0$.
    
    Para $i\in \{1, 2\}$, como $L_i$ é limite de $f$ em $p$, existe $\delta_i > 0$ tal que, para todo $x \in A\setminus\{p\}$, se $d(x, p) < \delta_i$, então $d(f(x), L_i) < R$.

    Seja $\delta = \min\{\delta_1, \delta_2\}$.
    Como $p$ é ponto de acumulação de $A$, existe $x \in A\setminus\{p\}$ tal que $d(x, p) < \delta$.
    Assim, $d(x, p)<\delta\leq \delta_1, \delta_2$, logo:

    \[
        d(f(x), L_1) < R \qquad\text{e}\qquad d(f(x), L_2) < R.
    \]
    Assim, da desigualdade triangular, temos que:
    \[
        d(L_1, L_2) \leq d(L_1, f(x)) + d(f(x), L_2) < R + R = 2R = d(L_1, L_2).
    \]
    Como $d(L_1, L_2) < d(L_1, L_2)$ é impossível, chegamos a uma contradição. Logo, $L_1 = L_2$.
\end{proof}


\subimport{}{limiteUnico-figura}

Com isso, podemos definir uma notação muito usual para se falar de limites.
\begin{definition}\index{limite!notação}
    Seja $A\subseteq M$ um conjunto, $f: A \to N$ uma função e $p \in A$ um ponto de acumulação de $A$.

    Então, caso exista, o único limite de $f$ em $p$ é denotado por
    
    \[
    \lim_{x\to p} f(x).
    \]
\end{definition}
A variável $x$, na definição acima, é uma variável muda.
No caso de funções de várias variáveis, é comum usar notações alternativas quando conveniente.
Por exemplo, se $f:\mathbb R^2 \to \mathbb R$ é uma função e $p=(a, b)$, é comum usar a notação
\[
    \lim_{(x, y)\to (a, b)} f(x, y).
\]

Para dimensões maiores, notações análogas são utilizadas.
\smallskip

A proposição abaixo conecta os conceitos de limite e de continuidade.
\begin{proposition}\index{continuidade!caracterização por limite}
    Sejam:
    \begin{itemize}
        \item $A\subseteq M$ um conjunto,
        \item $f: A \to N$ uma função, e
        \item $p \in A$ um ponto de acumulação de $A$.
    \end{itemize}

    São equivalentes:
    \begin{enumerate}[label=(\roman*)]
        \item $f$ é contínua em $p$;
        \item $f(p)$ é limite de $f$ em $p$.
    \end{enumerate}
\end{proposition}

\begin{proof}
    Suponha que $f$ é contínua em $p$.
    Veremos que $f(p)$ é limite de $f$ em $p$.

    Para tanto, fixe $\epsilon>0$.
    Devemos ver que existe $\delta>0$ tal que, para todo $x \in A\setminus\{p\}$, se $d(x, p) < \delta$, então $d'(f(x), f(p)) < \epsilon$.
    Ora, como $f$ é contínua em $p$, temos que tal $\delta$ existe.
    \smallskip
    
    Reciprocamente, suponha que $f$ possui limite em $p$ e que tal limite é $f(p)$.
    Devemos ver que $f$ é contínua em $p$.

    Para isso, fixamos $\epsilon>0$.
    Como $f(p)$ é limite de $f$ em $p$, existe $\delta>0$ tal que, para todo $x \in A\setminus\{p\}$, se $d(x, p) < \delta$, então $d'(f(x), f(p)) < \epsilon$.

    Veremos que tal $\delta$ funciona para a definição de continuidade.

    De fato, devemos ver que se $x \in A$ e $d(x, p) < \delta$, então $d'(f(x), f(p)) < \epsilon$.
    Isso decorre da escolha de $\delta$ caso $x \in A\setminus\{p\}$.
    Já se $x = p$, então $d'(f(x), f(p)) = d'(f(p), f(p)) = 0 < \epsilon$.
\end{proof}

Assim, temos o seguinte valioso teste para se decidir se uma função é contínua em um ponto.


O seguinte corolário sintetiza uma forma de se decidir se uma função é ou não é contínua.

\begin{corollary}
    Sejam:
    \begin{itemize}
        \item $A\subseteq M$ um conjunto,
        \item $f: A \to N$ uma função, e
        \item $p \in A$.
    \end{itemize}
    Então:
    \begin{enumerate}[label=(\roman*)]
        \item se $p$ é um ponto isolado de $A$, então $f$ é contínua em $p$;
        \item se $p$ é um ponto de acumulação de $A$ e $f$ não possui limite em $p$, então $f$ não é contínua em $p$;
        \item se $p$ é um ponto de acumulação de $A$ e $f$ possui limite em $p$ que é diferente de $f(p)$, então $f$ não é contínua em $p$;
        \item se $p$ é um ponto de acumulação de $A$ e $f(p)$ é limite de $f$ em $p$, então $f$ é contínua em $p$.
    \end{enumerate}
\end{corollary}


\begin{proposition}\index{limite!via funções coordenadas}Sejam
    \begin{itemize}
        \item $A\subseteq \mathbb R^n$ um conjunto,
        \item $f: A \to \mathbb R^m$ uma função, e
        \item $p \in A$ um ponto de acumulação de $A$.
        \item $L=(L_1, \dots, L_m)\in \mathbb R^m$.
    \end{itemize}

    Então:
    \[
        L \text{ é limite de } f \text{ em } p \text{ se, e somente se, } \forall i \in \{1, \dots, m\}\, L_i \text{ é limite de } f_i \text{ em } p.
    \]
\end{proposition}

\begin{proof}
    Primeiro, suponha que $L$ é limite de $f$ em $p$.
    Vejamos que cada $L_i$ é limite de $f_i$ em $p$.
    
    Para isso, fixamos um $i \in \{1, \dots, m\}$ e consideramos $\epsilon>0$.

    Como $L$ é limite de $f$ em $p$, existe $\delta>0$ tal que, para todo $x \in A\setminus\{p\}$, se $d(x, p) < \delta$, então $d(f(x), L) < \epsilon$.

    Veremos que o mesmo $\delta$ funciona para $f_i$.
    De fato, se $d(x, p) < \delta$, então temos:

    \begin{align*}
        |f_i(x) - L_i| &= \sqrt{(f_i(x) - L_i)^2} \\
        &= \sqrt{(f_1(x) - L_1)^2 + \cdots + (f_m(x) - L_m)^2} \\
        &= \|f(x) - L\| = d(f(x), L) < \epsilon.
    \end{align*}

    Portanto, $f_i$ é limite de $f_i$ em $p$.

    Reciprocamente, suponha que para todo $i$, $L_i$ é limite de $f_i$ em $p$.
    Fixamos $\epsilon>0$ e, para cada $i$, existe $\delta_i>0$ tal que, se $x \in A\setminus\{p\}$ e $d(x, p) < \delta_i$, então $|f_i(x) - L_i| < \epsilon$.

    Seja $\delta = \min\{\delta_1, \ldots, \delta_m\}$.
    Então, se $x \in A\setminus\{p\}$ e $d(x, p) < \delta$, temos que:

    \begin{align*}
        d(f(x), L) &= \sqrt{(f_1(x) - L_1)^2 + \cdots + (f_m(x) - L_m)^2} \\
        &< \sqrt{\epsilon^2+\dots+\epsilon^2}= \sqrt{m}\epsilon.
    \end{align*}

    Assim, se de partida, ao invés de $\epsilon$, tomarmos $\delta_1, \dots, \delta_n$ que funcione para $\epsilon'=\frac{\epsilon}{\sqrt{m}}$, concluiremos que para todo $x \in A\setminus\{p\}$ com $d(x, p) < \delta$, temos que $d(f(x), L) < \epsilon'$.
\end{proof}

\begin{proposition}Sejam
    \begin{itemize}
        \item $A\subseteq \mathbb R^n$ um conjunto,
        \item $B\subseteq A$,
        \item $p$ ponto de acumulação de $B$,
        \item $f: A \to \mathbb R^m$ uma função, e
        \item $L=\lim_{x\to p} f(x)$.
    \end{itemize}

    Então
    \[
    \lim_{x\to p} (f|_B)(x) = L.
    \]
\end{proposition}

\begin{proposition}Sejam
    \begin{itemize}
        \item $A\subseteq \mathbb R^n$ um conjunto,
        \item $B\subseteq A$,
        \item $p\in B$, e
        \item $f: A \to \mathbb R^m$ uma função.
    \end{itemize}

    Se $f$ é contínua em $p$, então $f|_B$ é contínua em $p$.
\end{proposition}


\begin{proposition}Sejam
    \begin{itemize}
        \item $A\subseteq \mathbb R^n$ um conjunto,
        \item $B\subseteq \mathbb R^m$ um conjunto,
        \item $p \in \mathbb R^n$ um ponto de acumulação de $A$,
        \item $f: A \to B$ uma função,
        \item $g: B \to \mathbb R^k$ uma função, e
        \item $L=\lim_{x\to p} f(x)$.
    \end{itemize}

    Se $L\in B$ e $g$ é contínua em $L$, então
    \[
    \lim_{x\to p} (g\circ f)(x) = g(L).
    \]
\end{proposition}

\begin{proposition}Sejam
    \begin{itemize}
        \item $A\subseteq \mathbb R^n$ um conjunto,
        \item $B\subseteq \mathbb R^m$ um conjunto,
        \item $p \in A$,
        \item $f: A \to B$ uma função,
        \item $g: B \to \mathbb R^k$ uma função.
    \end{itemize}

    Se $f$ é contínua em $p$ e $g$ é contínua em $f(p)$, então $g\circ f$ é contínua em $p$.

    Consequentemente, se $f$ e $g$ são contínuas, então $g\circ f$ é contínua.
\end{proposition}
